
%DIF LATEXDIFF DIFFERENCE FILE
%DIF DEL lecture5.tex             Thu Nov  6 02:44:42 2025
%DIF ADD diffs/lecture5_new.tex   Fri Sep 25 16:34:51 2026
\documentclass [a4paper,12pt]{paper}

\usepackage{mycommands}
\title{Lecture 16: Ray theory}
\author{David Al-Attar \\ Michaelmas Term, \the\year}
%DIF PREAMBLE EXTENSION ADDED BY LATEXDIFF
%DIF UNDERLINE PREAMBLE %DIF PREAMBLE
\RequirePackage[normalem]{ulem} %DIF PREAMBLE
\RequirePackage{color}\definecolor{RED}{rgb}{1,0,0}\definecolor{BLUE}{rgb}{0,0,1} %DIF PREAMBLE
\providecommand{\DIFadd}[1]{{\protect\color{blue}\uwave{#1}}} %DIF PREAMBLE
\providecommand{\DIFdel}[1]{{\protect\color{red}\sout{#1}}}                      %DIF PREAMBLE
%DIF SAFE PREAMBLE %DIF PREAMBLE
\providecommand{\DIFaddbegin}{} %DIF PREAMBLE
\providecommand{\DIFaddend}{} %DIF PREAMBLE
\providecommand{\DIFdelbegin}{} %DIF PREAMBLE
\providecommand{\DIFdelend}{} %DIF PREAMBLE
\providecommand{\DIFmodbegin}{} %DIF PREAMBLE
\providecommand{\DIFmodend}{} %DIF PREAMBLE
%DIF FLOATSAFE PREAMBLE %DIF PREAMBLE
\providecommand{\DIFaddFL}[1]{\DIFadd{#1}} %DIF PREAMBLE
\providecommand{\DIFdelFL}[1]{\DIFdel{#1}} %DIF PREAMBLE
\providecommand{\DIFaddbeginFL}{} %DIF PREAMBLE
\providecommand{\DIFaddendFL}{} %DIF PREAMBLE
\providecommand{\DIFdelbeginFL}{} %DIF PREAMBLE
\providecommand{\DIFdelendFL}{} %DIF PREAMBLE
%DIF COLORLISTINGS PREAMBLE %DIF PREAMBLE
\RequirePackage{listings} %DIF PREAMBLE
\RequirePackage{color} %DIF PREAMBLE
\lstdefinelanguage{DIFcode}{ %DIF PREAMBLE
%DIF DIFCODE_UNDERLINE %DIF PREAMBLE
  moredelim=[il][\color{red}\sout]{\%DIF\ <\ }, %DIF PREAMBLE
  moredelim=[il][\color{blue}\uwave]{\%DIF\ >\ } %DIF PREAMBLE
} %DIF PREAMBLE
\lstdefinestyle{DIFverbatimstyle}{ %DIF PREAMBLE
	language=DIFcode, %DIF PREAMBLE
	basicstyle=\ttfamily, %DIF PREAMBLE
	columns=fullflexible, %DIF PREAMBLE
	keepspaces=true %DIF PREAMBLE
} %DIF PREAMBLE
\lstnewenvironment{DIFverbatim}{\lstset{style=DIFverbatimstyle}}{} %DIF PREAMBLE
\lstnewenvironment{DIFverbatim*}{\lstset{style=DIFverbatimstyle,showspaces=true}}{} %DIF PREAMBLE
%DIF END PREAMBLE EXTENSION ADDED BY LATEXDIFF

\begin{document}

\maketitle
\DIFdelbegin %DIFDELCMD < 

%DIFDELCMD < %%%
\DIFdelend \section*{Outline and \DIFdelbegin \DIFdel{Motivation}\DIFdelend \DIFaddbegin \DIFadd{motivation}\DIFaddend }

In the previous lecture we examined the propagation of plane elastic waves. Here we consider what happens when the material 
parameters are not constant, 
developing a ray theory that provides a \textbf{good approximation in situations when the \DIFdelbegin \DIFdel{length-scale }\DIFdelend \DIFaddbegin \DIFadd{length scale }\DIFaddend of the 
heterogeneity is large relative to that of the waves}. Essentially, the waves always look locally 
planar, but as they propagate the wavefronts gradually deform, and the amplitudes can be changed by focusing 
and defocusing effects. This ray theory is related to elastodynamics in the same manner that 
geometric optics is related to Maxwell's equations, with both being \DIFdelbegin \DIFdel{"high frequency" }\DIFdelend \DIFaddbegin \DIFadd{``high frequency'' }\DIFaddend approximations in a 
suitable sense. There are also close links with what is known as semi-classical analysis, this 
being the subject in which the behaviour of quantum mechanical systems \DIFdelbegin \DIFdel{are }\DIFdelend \DIFaddbegin \DIFadd{is }\DIFaddend studied in the \DIFdelbegin \DIFdel{"classical limit" }\DIFdelend \DIFaddbegin \DIFadd{``classical limit'' }\DIFaddend $\hbar\rightarrow0$.

\section*{Relative length scales of heterogeneity}

We consider the linearised motion of an elastic body occupying a whole space whose material parameters are 
constant everywhere except a bounded region in which $\rho$ and $A_{ijkl}$ can vary smoothly. 
For definiteness, let $p_{i}^{0}$ be the slowness vector of \DIFdelbegin \DIFdel{elastic waves within }\DIFdelend \DIFaddbegin \DIFadd{a chosen plane wave in }\DIFaddend the homogeneous part of the body, and 
assume that the heterogeneous region is contained within the \DIFdelbegin \DIFdel{half-plane }\DIFdelend \DIFaddbegin \DIFadd{half-space }\DIFaddend $p_{i}^{0}x_{i}\ge0$. 
Within the homogeneous \DIFdelbegin \DIFdel{half plane }\DIFdelend \DIFaddbegin \DIFadd{half-space }\DIFaddend $p_{i}^{0}x_{i}<0$ we know that plane wave solutions are possible of the form
\DIFdelbegin \begin{displaymath}
\DIFdel{u_{i}(x,t)=f(t-p_{i}^{0}x_{i})a_{i}^{0}
}\end{displaymath}%DIFAUXCMD
\DIFdelend \DIFaddbegin \begin{equation}
\DIFadd{u_{i}(\mathbf{x},t)=f(t-p_{j}^{0}x_{j})a_{i}^{0},
\label{eq:1}
}\end{equation}\DIFaddend 
and our aim is to extend this solution suitably into the heterogeneous region. In general there exists no exact
closed-form solution to this problem, and so it must be studied either through numerical calculations 
or using approximate methods. While numerical calculations are of huge practical value, they provide 
relatively little physical insight, and so we shall focus on physically motivated approximate methods.


\begin{figure}[ht]
\centering
\DIFdelbeginFL %DIFDELCMD < \includegraphics[width=0.8\textwidth]{figures/L5F1.png}
%DIFDELCMD < %%%
\DIFdelendFL \DIFaddbeginFL \includegraphics[width=0.8\textwidth]{figures/L5F1.pdf}
\DIFaddendFL \caption{\DIFdelbeginFL \DIFdelFL{Scattering }\DIFdelendFL \DIFaddbeginFL \DIFaddFL{Snapshot from a numerical simulation of the scattering }\DIFaddendFL of a plane \DIFdelbeginFL \DIFdelFL{p-wave }\DIFdelendFL \DIFaddbeginFL \DIFaddFL{P-wave }\DIFaddendFL in a homogeneous background model due to its \DIFdelbeginFL \DIFdelFL{inter-
action }\DIFdelendFL \DIFaddbeginFL \DIFaddFL{interaction 
}\DIFaddendFL with a heterogeneity having a typical \DIFdelbeginFL \DIFdelFL{length-scale }\DIFdelendFL \DIFaddbeginFL \DIFaddFL{length scale }\DIFaddendFL approximately equal to
the \DIFdelbeginFL \DIFdelFL{length-scale }\DIFdelendFL \DIFaddbeginFL \DIFaddFL{length scale }\DIFaddendFL of the incident wave. The \DIFaddbeginFL \DIFaddFL{colours show the magnitude of the displacement, the }\DIFaddendFL arrows show the particle motion, and \DIFdelbeginFL \DIFdelFL{we see that }\DIFdelendFL the \DIFaddbeginFL \DIFaddFL{dashed line marks the extent of the heterogeneity; the colour scale is saturated so as to show the weak }\DIFaddendFL scattered \DIFdelbeginFL \DIFdelFL{wave }\DIFdelendFL field\DIFaddbeginFL \DIFaddFL{. The box }\DIFaddendFL is \DIFdelbeginFL \DIFdelFL{comprised }\DIFdelendFL \DIFaddbeginFL \DIFaddFL{about 13 wavelengths }\DIFaddendFL of \DIFaddbeginFL \DIFaddFL{the incident wave wide and 6 high. We see that the scattered wavefield comprises }\DIFaddendFL both P-waves and S-waves.}
\DIFdelbeginFL %DIFDELCMD < \label{fig:fig1}
%DIFDELCMD < %%%
\DIFdelendFL \DIFaddbeginFL \label{fig:1}
\DIFaddendFL \end{figure}

In broad terms there are three possible approaches for studying this problem, with the appropriate choice being determined
by the relative length scales of the incident wave and the heterogeneous structure. Let $\lambda$ 
\DIFaddbegin \DIFadd{(not to be confused with the Lam\'{e} parameter) }\DIFaddend denote the \textbf{dominant wavelength for the plane wave}, and denote by $L$ the \textbf{smallest length scale over which 
the material parameters vary to an appreciable extent}. If their ratio satisfies
\DIFdelbegin \begin{displaymath}
\DIFdel{\frac{L}{\lambda}\approx1
}\end{displaymath}%DIFAUXCMD
\DIFdelend \DIFaddbegin \begin{equation}
\DIFadd{\frac{L}{\lambda}\approx1,
\label{eq:2}
}\end{equation}\DIFaddend 
then we expect the incident wave to be \textbf{scattered by the heterogeneity} \DIFaddbegin \DIFadd{(see Fig.~\ref{fig:1})}\DIFaddend . There is not sufficient time in this course to 
discuss scattering of seismic waves further, though it is an important physical process within the Earth, 
and can be analysed using the \textbf{Born approximation} that you may know from quantum mechanics. 

In this lecture we shall focus on situations when
\DIFdelbegin \begin{displaymath}
\DIFdel{\frac{L}{\lambda}\gg1
}\end{displaymath}%DIFAUXCMD
\DIFdelend \DIFaddbegin \begin{equation}
\DIFadd{\frac{L}{\lambda}\gg1,
\label{eq:3}
}\end{equation}\DIFaddend 
which is to say the length scale of the heterogeneity is much greater than the length scale of the incident wave. Here we 
do not expect to see significant scattering of the incident wave, but instead a smooth evolution of the 
wavefront as it propagates through the heterogeneity, along with associated amplitude variations due to focusing effects \DIFaddbegin \DIFadd{(see Fig.~\ref{fig:2})}\DIFaddend . This 
is the domain of \textbf{geometric ray theory}.


Finally, we might ask what happens in the case that
\DIFdelbegin \begin{displaymath}
\DIFdel{\frac{L}{\lambda}\ll1
}\end{displaymath}%DIFAUXCMD
\DIFdelend \DIFaddbegin \begin{equation}
\DIFadd{\frac{L}{\lambda}\ll1.
\label{eq:4}
}\end{equation}\DIFaddend 
This is actually very applicable to the Earth, as rocks are made up of small grains of different minerals that can have 
quite distinct physical properties, and yet observed seismic waves have wavelengths on the order of kilometres 
to thousands of kilometres. We also lack the time to enter into this topic in any detail, but it can be shown 
that such waves ``see'' a suitably averaged version of the body that appears smooth on the length scale of the incident waves.
The appropriate smoothed medium can be determined using the methods of \textbf{homogenisation theory}. Again, there
is sadly not time within this course to discuss this topic further.


\begin{figure}[ht]
\centering
\DIFdelbeginFL %DIFDELCMD < \includegraphics[width=0.8\textwidth]{figures/L5F2.png}
%DIFDELCMD < %%%
\DIFdelendFL \DIFaddbeginFL \includegraphics[width=0.8\textwidth]{figures/L5F2.pdf}
\DIFaddendFL \caption{\DIFdelbeginFL \DIFdelFL{Passage }\DIFdelendFL \DIFaddbeginFL \DIFaddFL{As for Fig.~\ref{fig:1}, but for the passage }\DIFaddendFL of a plane \DIFdelbeginFL \DIFdelFL{p-wave }\DIFdelendFL \DIFaddbeginFL \DIFaddFL{P-wave }\DIFaddendFL through a heterogeneous region having \DIFdelbeginFL \DIFdelFL{length-scale }\DIFdelendFL \DIFaddbeginFL \DIFaddFL{length scale }\DIFaddendFL much
larger than that of the incident wave. Here we see minimal scattering, along with
the smooth distortion of the \DIFdelbeginFL \DIFdelFL{wave front }\DIFdelendFL \DIFaddbeginFL \DIFaddFL{wavefront }\DIFaddendFL and associated amplitude and polarisation
variations. This is a situation to which ray theory can be usefully applied.}
\DIFdelbeginFL %DIFDELCMD < \label{fig:fig2}
%DIFDELCMD < %%%
\DIFdelendFL \DIFaddbeginFL \label{fig:2}
\DIFaddendFL \end{figure}

\section*{The ray series ansatz}

As noted above, we shall focus on situations where the length scale of the heterogeneity is long 
relative to the length scale of the incident plane wave. As a first guess, we might expect that the continuation of 
the plane wave into the heterogeneous region can be approximated in the following manner
\DIFdelbegin \begin{displaymath}
\DIFdel{u_{i}(x,t)=f}[\DIFdel{t-T(x)}]\DIFdel{a_{i}(x)
}\end{displaymath}%DIFAUXCMD
\DIFdelend \DIFaddbegin \begin{equation}
\DIFadd{u_{i}(\mathbf{x},t)=f}[\DIFadd{t-T(\mathbf{x})}]\DIFadd{a_{i}(\mathbf{x}),
\label{eq:5}
}\end{equation}\DIFaddend 
with this expression taking account of the following physical ideas:
\begin{enumerate}
    \item The linear phase term $t-p_{i}^{0}x_{i}$ has been replaced with the non-linear term \DIFdelbegin \DIFdel{$t-T(x)$}\DIFdelend \DIFaddbegin \DIFadd{$t-T(\mathbf{x})$}\DIFaddend . This allows the 
    wavefronts in the heterogeneous region to curve as they \DIFdelbegin \DIFdel{propagates }\DIFdelend \DIFaddbegin \DIFadd{propagate }\DIFaddend through regions with different wave speeds;
    \item The amplitude and polarisation term $a_{i}^{0}$ has been replaced with a function of position \DIFdelbegin \DIFdel{$a_{i}(x)$ }\DIFdelend \DIFaddbegin \DIFadd{$a_{i}(\mathbf{x})$ }\DIFaddend which 
    allows for (i) focusing and defocusing effects during propagation and (ii) changes in the particle motion associated 
    with both wavefront curvature and variations in anisotropic properties.
\end{enumerate}
We will call the function $T$ the \textbf{travel time}, and note that the wavefronts of eq.\DIFdelbegin \DIFdel{(5}\DIFdelend \DIFaddbegin \DIFadd{~(\ref{eq:5}}\DIFaddend ) are, by definition\DIFdelbegin \DIFdel{level surfaces
}\begin{displaymath}
\DIFdel{T(x)=\text{constant}.
}\end{displaymath}%DIFAUXCMD
\DIFdelend \DIFaddbegin \DIFadd{, the level surfaces
}\begin{equation}
\DIFadd{T(\mathbf{x})=\text{constant}.
\label{eq:6}
}\end{equation}\DIFaddend 
It will be useful to define the \textbf{local slowness vector} in terms of this travel time by setting
\DIFdelbegin \begin{displaymath}
\DIFdel{p_{i}=\frac{\partial T}{\partial x_{i}}
}\end{displaymath}%DIFAUXCMD
\DIFdelend \DIFaddbegin \begin{equation}
\DIFadd{p_{i}=\frac{\partial T}{\partial x_{i}},
\label{eq:7}
}\end{equation}\DIFaddend 
which is everywhere orthogonal to the wavefronts. Note that the local slowness vector is equal 
to the slowness vector within homogeneous regions where the travel time depends linearly on \DIFdelbegin \DIFdel{$\mathbf{p}$}\DIFdelend \DIFaddbegin \DIFadd{$\mathbf{x}$}\DIFaddend .

To proceed, it is useful to \DIFdelbegin \DIFdel{decompose displacement vector fields into an integralover harmonic plane waves
}\begin{displaymath}
\DIFdel{u_{i}(x,t)=\frac{1}{2\pi}\int_{\mathbb{R}}\tilde{f}(\omega)e^{\ii[t-T(x)]}a_{i}(x)\dd\omega,
}\end{displaymath}%DIFAUXCMD
\DIFdelend \DIFaddbegin \DIFadd{write the waveform $f$ as a Fourier integral, so that eq.~(\ref{eq:5}) becomes
}\begin{equation}
\DIFadd{u_{i}(\mathbf{x},t)=\frac{1}{2\pi}\int_{\mathbb{R}}\tilde{f}(\omega)\ee^{\ii\omega[t-T(\mathbf{x})]}a_{i}(\mathbf{x})\dd\omega,
\label{eq:8}
}\end{equation}\DIFaddend 
\DIFaddbegin \DIFadd{where $\tilde{f}$ denotes the Fourier transform of $f$. }\DIFaddend It is clear that eq.\DIFdelbegin \DIFdel{(5}\DIFdelend \DIFaddbegin \DIFadd{~(\ref{eq:5}}\DIFaddend ) is the temporal convolution of the singular plane wave
\DIFdelbegin \begin{displaymath}
\DIFdel{u_{i}(x,t)=\delta[t-T(x)]a_{i}(x)
}\end{displaymath}%DIFAUXCMD
\DIFdelend \DIFaddbegin \begin{equation}
\DIFadd{u_{i}(\mathbf{x},t)=\delta[t-T(\mathbf{x})]a_{i}(\mathbf{x})
\label{eq:9}
}\end{equation}\DIFaddend 
with the waveform \DIFdelbegin \DIFdel{f}\DIFdelend \DIFaddbegin \DIFadd{$f$}\DIFaddend . In what follows, it will be convenient to focus on such singular waves, and from now on 
we set $\tilde{f}(\omega)=1$ or equivalently $f(t)=\delta(t)$.

We could try to use the ansatz in eq.\DIFdelbegin \DIFdel{(9}\DIFdelend \DIFaddbegin \DIFadd{~(\ref{eq:9}}\DIFaddend ) as a 
starting point for the development of ray theory, but experience has shown that it is necessary to consider 
instead the following generalisation
\DIFdelbegin \begin{displaymath}
\DIFdel{u_{i}(x,t)=\frac{1}{2\pi}\int_{\mathbb{R}}\sum_{n=0}^{\infty}\frac{a_{i}^{n}(x)}{(\ii\omega)^{n}}e^{\ii\omega[t-T(x)]}\dd\omega,
}\end{displaymath}%DIFAUXCMD
\DIFdelend \DIFaddbegin \begin{equation}
\DIFadd{u_{i}(\mathbf{x},t)=\frac{1}{2\pi}\int_{\mathbb{R}}\sum_{n=0}^{\infty}\frac{a_{i}^{n}(\mathbf{x})}{(\ii\omega)^{n}}\ee^{\ii\omega[t-T(\mathbf{x})]}\dd\omega,
\label{eq:10}
}\end{equation}\DIFaddend 
where we note that the $n$ in $a_{i}^{n}$ is a superscript and not an exponent. This expression is known as the \textbf{elastodynamic 
ray series}. The $n=0$ term within this expression is equivalent to \DIFdelbegin \DIFdel{the eq.(9}\DIFdelend \DIFaddbegin \DIFadd{eq.~(\ref{eq:9}}\DIFaddend ). The physical meaning or need for the additional terms 
within the expansion is less clear, and indeed, these terms are typically ignored within practical applications. But 
as an initial comment, we note that each successive term within the series is \DIFdelbegin \DIFdel{divided by a }\DIFdelend \DIFaddbegin \DIFadd{multiplied by a further }\DIFaddend factor of $1/(\ii\omega)$, and within the Fourier domain 
this is equivalent to a temporal integration. Thus the terms within this series are getting progressively smoother and smoother.

\subsection*{Substitution into the equations of motion}
We recall that the linearised equations of motion for an elastic body can be written
\DIFdelbegin \begin{displaymath}
\DIFdel{\rho\frac{\partial^{2}u_{i}}{\partial t^{2}}-\frac{\partial}{\partial x_{j}}\left(A_{ijkl}\frac{\partial u_{k}}{\partial x_{l}}\right)=0,
}\end{displaymath}%DIFAUXCMD
\DIFdelend \DIFaddbegin \begin{equation}
\DIFadd{\rho\frac{\partial^{2}u_{i}}{\partial t^{2}}-\frac{\partial}{\partial x_{j}}\left(A_{ijkl}\frac{\partial u_{k}}{\partial x_{l}}\right)=0,
\label{eq:11}
}\end{equation}\DIFaddend 
where we \DIFdelbegin \DIFdel{are not here allowing }\DIFdelend \DIFaddbegin \DIFadd{do not include }\DIFaddend a body force associated with a seismic source. Within the temporal Fourier domain this equation becomes
\DIFdelbegin \begin{displaymath}
\DIFdel{-\omega^{2}\rho u_{i}-\frac{\partial}{\partial x_{j}}\left(A_{ijkl}\frac{\partial u_{k}}{\partial x_{l}}\right)=0
}\end{displaymath}%DIFAUXCMD
\DIFdelend \DIFaddbegin \begin{equation}
\DIFadd{-\omega^{2}\rho u_{i}-\frac{\partial}{\partial x_{j}}\left(A_{ijkl}\frac{\partial u_{k}}{\partial x_{l}}\right)=0,
\label{eq:12}
}\end{equation}\DIFaddend 
where for simplicity we use the same symbol for the displacement vector and its Fourier transform. From eq.\DIFdelbegin \DIFdel{(10}\DIFdelend \DIFaddbegin \DIFadd{~(\ref{eq:10}}\DIFaddend ) we see that the 
ray series in the Fourier domain is simply
\DIFdelbegin \begin{displaymath}
\DIFdel{u_{i}(x,\omega)=\sum_{n=0}^{\infty}\frac{a_{i}^{n}(x)}{(\ii\omega)^{n}}e^{-\ii\omega T(x)}.
}\end{displaymath}%DIFAUXCMD
\DIFdelend \DIFaddbegin \begin{equation}
\DIFadd{u_{i}(\mathbf{x},\omega)=\sum_{n=0}^{\infty}\frac{a_{i}^{n}(\mathbf{x})}{(\ii\omega)^{n}}\ee^{-\ii\omega T(\mathbf{x})}.
\label{eq:13}
}\end{equation}\DIFaddend 
From \DIFdelbegin \DIFdel{eq.}\DIFdelend \DIFaddbegin \DIFadd{eqs.~(\ref{eq:7}) and }\DIFaddend (\DIFdelbegin \DIFdel{7) and eq.(13}\DIFdelend \DIFaddbegin \DIFadd{\ref{eq:13}}\DIFaddend ) we find that
\DIFdelbegin \begin{eqnarray*}
\DIFdel{\frac{\partial u_{k}}{\partial x_{l}}}&\DIFdel{=\sum_{n=0}^{\infty}\frac{1}{(\ii\omega)^{n}}\left(\frac{\partial a_{k}^{n}}{\partial x_{l}}
-\ii\omega p_{l}a_{k}^{n}\right)e^{-\ii\omega T}, }\\
\DIFdel{A_{ijkl}\frac{\partial u_{k}}{\partial x_{l}}}&\DIFdel{=\sum_{n=0}^{\infty}\frac{1}{(\ii\omega)^{n}}A_{ijkl}\left(\frac{\partial a_{k}^{n}}{\partial x_{l}}-\ii\omega p_{l}a_{k}^{n}\right)e^{-\ii\omega T}.
}\end{eqnarray*}%DIFAUXCMD
\DIFdelend \DIFaddbegin \begin{align}
\DIFadd{\frac{\partial u_{k}}{\partial x_{l}}}&\DIFadd{=\sum_{n=0}^{\infty}\frac{1}{(\ii\omega)^{n}}\left(\frac{\partial a_{k}^{n}}{\partial x_{l}}
-\ii\omega p_{l}a_{k}^{n}\right)\ee^{-\ii\omega T}, \label{eq:14}}\\
\DIFadd{A_{ijkl}\frac{\partial u_{k}}{\partial x_{l}}}&\DIFadd{=\sum_{n=0}^{\infty}\frac{1}{(\ii\omega)^{n}}A_{ijkl}\left(\frac{\partial a_{k}^{n}}{\partial x_{l}}-\ii\omega p_{l}a_{k}^{n}\right)\ee^{-\ii\omega T}. \label{eq:15}
}\end{align}\DIFaddend 
Taking the divergence of the second expression, we obtain
\DIFdelbegin \begin{eqnarray*}
\DIFdel{\frac{\partial}{\partial x_{j}}\left(A_{ijkl}\frac{\partial u_{k}}{\partial x_{l}}\right)=\sum_{n=0}^{\infty}\frac{1}{(\ii\omega)^{n}}\left[\frac{\partial}{\partial x_{j}}\left(A_{ijkl}\frac{\partial a_{k}^{n}}{\partial x_{l}}\right)
-\ii\omega\frac{\partial}{\partial x_{j}}(A_{ijkl}p_{l}a_{k}^{n}) \right. }\\
\DIFdel{\left. -\ii\omega A_{ijkl}p_{j}\frac{\partial a_{k}^{n}}{\partial x_{l}}+(\ii\omega)^{2}A_{ijkl}p_{j}p_{l}a_{k}^{n}\right]e^{-\ii\omega T}.
}\end{eqnarray*}%DIFAUXCMD
\DIFdelend \DIFaddbegin \begin{multline}
\DIFadd{\frac{\partial}{\partial x_{j}}\left(A_{ijkl}\frac{\partial u_{k}}{\partial x_{l}}\right)=\sum_{n=0}^{\infty}\frac{1}{(\ii\omega)^{n}}\left[\frac{\partial}{\partial x_{j}}\left(A_{ijkl}\frac{\partial a_{k}^{n}}{\partial x_{l}}\right)
-\ii\omega\frac{\partial}{\partial x_{j}}(A_{ijkl}p_{l}a_{k}^{n}) \right. }\\
\DIFadd{\left. -\ii\omega A_{ijkl}p_{j}\frac{\partial a_{k}^{n}}{\partial x_{l}}+(\ii\omega)^{2}A_{ijkl}p_{j}p_{l}a_{k}^{n}\right]\ee^{-\ii\omega T},
\label{eq:16}
}\end{multline}\DIFaddend 
and substituting into the frequency domain equations of motion in eq.\DIFdelbegin \DIFdel{(12}\DIFdelend \DIFaddbegin \DIFadd{~(\ref{eq:12}}\DIFaddend ) we arrive at
\DIFdelbegin \begin{eqnarray*}
\DIFdel{\sum_{n=0}^{\infty}\frac{1}{(\ii\omega)^{n}}\left[\frac{\partial}{\partial x_{j}}\left(A_{ijkl}\frac{\partial a_{k}^{n}}{\partial x_{l}}\right)-\ii\omega\frac{\partial}{\partial x_{j}}(A_{ijkl}p_{l}a_{k}^{n}) \right. }\\
\DIFdel{\left. -\ii\omega A_{ijkl}p_{j}\frac{\partial a_{k}^{n}}{\partial x_{l}}+(\ii\omega)^{2}A_{ijkl}p_{j}p_{l}a_{k}^{n}-(\ii\omega)^{2}\rho a_{i}^{n}\right]e^{-\ii\omega T}=0 .
}\end{eqnarray*}%DIFAUXCMD
\DIFdelend \DIFaddbegin \begin{multline}
\DIFadd{\sum_{n=0}^{\infty}\frac{1}{(\ii\omega)^{n}}\left[\frac{\partial}{\partial x_{j}}\left(A_{ijkl}\frac{\partial a_{k}^{n}}{\partial x_{l}}\right)-\ii\omega\frac{\partial}{\partial x_{j}}(A_{ijkl}p_{l}a_{k}^{n}) \right. }\\
\DIFadd{\left. -\ii\omega A_{ijkl}p_{j}\frac{\partial a_{k}^{n}}{\partial x_{l}}+(\ii\omega)^{2}A_{ijkl}p_{j}p_{l}a_{k}^{n}-(\ii\omega)^{2}\rho a_{i}^{n}\right]\ee^{-\ii\omega T}=0.
\label{eq:17}
}\end{multline}\DIFaddend 

The above equation must hold for all frequencies, and this can only be true if the coefficient of each power of $\ii\omega$ is separately equal to zero. From the coefficients of the highest-order term $(\ii\omega)^{2}$ we obtain
\DIFdelbegin \begin{displaymath}
[\DIFdel{A_{ijkl}p_{j}p_{l}-\rho\delta_{ik}}]\DIFdel{a_{k}^{0}=0,
}\end{displaymath}%DIFAUXCMD
\DIFdelend \DIFaddbegin \begin{equation}
[\DIFadd{A_{ijkl}p_{j}p_{l}-\rho\delta_{ik}}]\DIFadd{a_{k}^{0}=0,
\label{eq:18}
}\end{equation}\DIFaddend 
while those for the $\ii\omega$ term take the form
\DIFdelbegin \begin{displaymath}
[\DIFdel{A_{ijkl}p_{j}p_{l}-\rho\delta_{ik}}]\DIFdel{a_{k}^{1}=\frac{\partial}{\partial x_{j}}(A_{ijkl}p_{l}a_{k}^{0})+A_{ijkl}p_{j}\frac{\partial a_{k}^{0}}{\partial x_{l}},
}\end{displaymath}%DIFAUXCMD
\DIFdelend \DIFaddbegin \begin{equation}
[\DIFadd{A_{ijkl}p_{j}p_{l}-\rho\delta_{ik}}]\DIFadd{a_{k}^{1}=\frac{\partial}{\partial x_{j}}(A_{ijkl}p_{l}a_{k}^{0})+A_{ijkl}p_{j}\frac{\partial a_{k}^{0}}{\partial x_{l}},
\label{eq:19}
}\end{equation}\DIFaddend 
and then for a general term with $n\ge2$
\DIFdelbegin \begin{displaymath}
[\DIFdel{A_{ijkl}p_{j}p_{l}-\rho\delta_{ik}}]\DIFdel{a_{k}^{n}=\frac{\partial}{\partial x_{j}}(A_{ijkl}p_{l}a_{k}^{n-1})+A_{ijkl}p_{j}\frac{\partial a_{k}^{n-1}}{\partial x_{l}}-\frac{\partial}{\partial x_{j}}(A_{ijkl}\frac{\partial a_{k}^{n-2}}{\partial x_{l}})
}\end{displaymath}%DIFAUXCMD
\DIFdelend \DIFaddbegin \begin{equation}
[\DIFadd{A_{ijkl}p_{j}p_{l}-\rho\delta_{ik}}]\DIFadd{a_{k}^{n}=\frac{\partial}{\partial x_{j}}(A_{ijkl}p_{l}a_{k}^{n-1})+A_{ijkl}p_{j}\frac{\partial a_{k}^{n-1}}{\partial x_{l}}-\frac{\partial}{\partial x_{j}}\left(A_{ijkl}\frac{\partial a_{k}^{n-2}}{\partial x_{l}}\right).
\label{eq:20}
}\end{equation}\DIFaddend 
The first of these equations depends on only the local slowness vector and the zeroth-order polarisation vector. We will see shortly that 
this equation can be solved for the travel time $T$ and also gives the orientation of \DIFdelbegin \DIFdel{$a^{0}(x)$ }\DIFdelend \DIFaddbegin \DIFadd{$\mathbf{a}^{0}(\mathbf{x})$ }\DIFaddend at each point. 
Having done this, solution of eq.\DIFdelbegin \DIFdel{(19}\DIFdelend \DIFaddbegin \DIFadd{~(\ref{eq:19}}\DIFaddend ) gives the amplitude of the zeroth-order term, while eq.\DIFdelbegin \DIFdel{(20}\DIFdelend \DIFaddbegin \DIFadd{~(\ref{eq:20}}\DIFaddend ) can be used to successively 
determine the higher-order polarisation vectors. We will not, however, discuss the solution of these higher-order equations
 within this course.

At this stage it is reasonable to ask where exactly we have made an approximation\DIFdelbegin \DIFdel{? }\DIFdelend \DIFaddbegin \DIFadd{. }\DIFaddend Assuming that the above set of equations 
can be solved for $T$ and all the \DIFdelbegin \DIFdel{$a^{n}$ }\DIFdelend \DIFaddbegin \DIFadd{$\mathbf{a}^{n}$ }\DIFaddend then it would seem that we have obtained an exact solution to the equations of motion. 
For this to be a valid conclusion, however, we must then think carefully about the sense in which the ray series converges. 
Though we will not enter into this topic in detail, the key point is \textbf{the ray series is not convergent in any usual sense}. 
Indeed, the problem is at low frequencies, with the series being very badly behaved at $\omega=0$. To force this series to converge, 
we could zero-out the solution in some neighbourhood of $\omega=0$, but by doing this we would be left with an approximate 
solution to the equations of motion. Because, however, the error lies only within a bounded range of frequencies it must be a smooth 
function. Building on such ideas, it can be shown that the ray theoretic solution \textbf{exactly predicts the 
 singular part of the wavefield}, but that there is a \textbf{smooth error term whose size need not be small}\DIFdelbegin \footnote{\DIFdel{These ideas are
formalised within the theory of }\textbf{\DIFdel{Fourier integral operators}} %DIFAUXCMD
\DIFdel{developed in the 1970s by Lars Hormander and others.}}%DIFAUXCMD
\addtocounter{footnote}{-1}%DIFAUXCMD
\DIFdel{.}\DIFdelend \DIFaddbegin \DIFadd{.}\footnote{\DIFadd{These ideas are
formalised within the theory of }\textbf{\DIFadd{Fourier integral operators}} \DIFadd{developed in the 1970s by Lars H\"{o}rmander and others (H\"{o}rmander 1971).}}
\DIFaddend 

\section*{The eikonal equation}

Within this section, we show that eq.\DIFdelbegin \DIFdel{(18}\DIFdelend \DIFaddbegin \DIFadd{~(\ref{eq:18}}\DIFaddend ) can be solved for the travel time $T$ using the so-called \textbf{method of characteristics}. 
This is a useful method that can be applied to both linear and non-linear first-order partial differential equations, 
and reduces their solution to the integration of systems of ordinary differential equations. In fact, it is through 
solution of these ordinary differential equations that rays will enter into the theory in a natural manner.

Before applying this technique, we must first reduce eq.\DIFdelbegin \DIFdel{(18) into }\DIFdelend \DIFaddbegin \DIFadd{~(\ref{eq:18}) to }\DIFaddend a single partial differential equation for the travel time. 
To do this, it will be useful to define the \textbf{local Christoffel \DIFdelbegin \DIFdel{operator}\DIFdelend \DIFaddbegin \DIFadd{matrix}\DIFaddend } through
\DIFdelbegin \begin{displaymath}
\DIFdel{\Gamma_{ik}(x,\mathbf{p})=\frac{1}{\rho(\mathbf{x})}A_{ijkl}(\mathbf{x})p_{j}p_{l}.
}\end{displaymath}%DIFAUXCMD
\DIFdelend \DIFaddbegin \begin{equation}
\DIFadd{\Gamma_{ik}(\mathbf{x},\mathbf{p})=\frac{1}{\rho(\mathbf{x})}A_{ijkl}(\mathbf{x})p_{j}p_{l},
\label{eq:21}
}\end{equation}\DIFaddend 
for all \DIFdelbegin \DIFdel{$x\in\mathbb{R}^{3}$ and vectors $\mathbf{a}$ }\DIFdelend \DIFaddbegin \DIFadd{$\mathbf{x}\in\mathbb{R}^{3}$ }\DIFaddend and $\mathbf{p}$. We can then write eq.\DIFdelbegin \DIFdel{(18}\DIFdelend \DIFaddbegin \DIFadd{~(\ref{eq:18}}\DIFaddend ) in the form
\DIFdelbegin \begin{displaymath}
\DIFdel{\{\bm{\Gamma}[\mathbf{x},\mathbf{p}(\mathbf{x})]-\mathbf{I}\}\mathbf{a}^{0}(\mathbf{x})=\mathbf{0},
}\end{displaymath}%DIFAUXCMD
\DIFdelend \DIFaddbegin \begin{equation}
\DIFadd{\{\bm{\Gamma}[\mathbf{x},\mathbf{p}(\mathbf{x})]-\mathbf{I}\}\mathbf{a}^{0}(\mathbf{x})=\mathbf{0},
\label{eq:22}
}\end{equation}\DIFaddend 
which closely resembles the Christoffel equation we encountered in the study of elastic plane waves. Indeed, 
the physical intuition underlying ray theory is that locally the solutions behave like plane waves.
 At each point, $\mathbf{x}$, we know that eq.\DIFdelbegin \DIFdel{(22}\DIFdelend \DIFaddbegin \DIFadd{~(\ref{eq:22}}\DIFaddend ) may be satisfied by a non-trivial polarisation vector 
 $\mathbf{a}^{0}(\mathbf{x})$ if and only if the local slowness vector $\mathbf{p}(\mathbf{x})$ is a solution of
\DIFdelbegin \begin{displaymath}
\DIFdel{\det\{\bm{\Gamma}[\mathbf{x},\mathbf{p}(\mathbf{x})]-\mathbf{I}\}=0,
}\end{displaymath}%DIFAUXCMD
\DIFdelend \DIFaddbegin \begin{equation}
\DIFadd{\det\{\bm{\Gamma}[\mathbf{x},\mathbf{p}(\mathbf{x})]-\mathbf{I}\}=0,
\label{eq:23}
}\end{equation}\DIFaddend 
which is known in the context of ray theory as an \textbf{eikonal equation}. Eq.\DIFdelbegin \DIFdel{(23}\DIFdelend \DIFaddbegin \DIFadd{~(\ref{eq:23}}\DIFaddend ) can be interpreted as requiring
 that at each $\mathbf{x}$, the local slowness vector $\mathbf{p}(\mathbf{x})$ lies on the \textbf{local slowness surface}.


Having determined the appropriate local slowness vector $\mathbf{p}(\mathbf{x})$, eq.\DIFdelbegin \DIFdel{(22}\DIFdelend \DIFaddbegin \DIFadd{~(\ref{eq:22}}\DIFaddend ) then implies that $\mathbf{a}^{0}(\mathbf{x})$ lies 
within the eigenspace of $\bm{\Gamma}[\mathbf{x},\hat{\mathbf{p}}(\mathbf{x})]$ with eigenvalue equal to the square of the corresponding 
local phase speed. For simplicity, we shall only consider materials whose local slowness surfaces are everywhere non-degenerate, and so we can write
\DIFdelbegin \begin{displaymath}
\DIFdel{\mathbf{a}^{0}(\mathbf{x})=A_{0}(\mathbf{x})\hat{\mathbf{a}}^{0}(\mathbf{x}),
}\end{displaymath}%DIFAUXCMD
\DIFdelend \DIFaddbegin \begin{equation}
\DIFadd{\mathbf{a}^{0}(\mathbf{x})=A_{0}(\mathbf{x})\hat{\mathbf{a}}^{0}(\mathbf{x}),
\label{eq:24}
}\end{equation}\DIFaddend 
where $A_{0}(\mathbf{x})$ \DIFaddbegin \DIFadd{is }\DIFaddend a scalar amplitude, and $\hat{\mathbf{a}}^{0}(\mathbf{x})$ is the appropriate normalised eigenvector which is 
determined uniquely up to its sign. \DIFdelbegin \DIFdel{We have assumed that the local slowness surfaces within the body are everywhere non-degenerate, and can therefore, }\DIFdelend \DIFaddbegin \DIFadd{In an isotropic medium the two S-wave sheets of the slowness surface coincide, and so this assumption 
fails; the reduction of the eikonal equation described below remains unambiguous in that case, but the polarisation of the S-wave is no longer 
determined uniquely by eq.~(\ref{eq:22}). Under this non-degeneracy assumption we can therefore }\DIFaddend always write the eikonal equation unambiguously in the form
\DIFdelbegin \begin{displaymath}
\DIFdel{\{||\mathbf{p}(\mathbf{x})||^{2}c_{1}[\mathbf{x},\hat{\mathbf{p}}(\mathbf{x})]^{2}-1\}\{||\mathbf{p}(\mathbf{x})||^{2}c_{2}[\mathbf{x},\hat{\mathbf{p}}(\mathbf{x})]^{2}-1\}\{||p(\mathbf{x})||^{2}c_{3}[\mathbf{x},\hat{\mathbf{p}}(\mathbf{x})]^{2}-1\}=0
}\end{displaymath}%DIFAUXCMD
\DIFdelend \DIFaddbegin \begin{equation}
\DIFadd{\{\|\mathbf{p}(\mathbf{x})\|^{2}c_{1}[\mathbf{x},\hat{\mathbf{p}}(\mathbf{x})]^{2}-1\}\{\|\mathbf{p}(\mathbf{x})\|^{2}c_{2}[\mathbf{x},\hat{\mathbf{p}}(\mathbf{x})]^{2}-1\}\{\|\mathbf{p}(\mathbf{x})\|^{2}c_{3}[\mathbf{x},\hat{\mathbf{p}}(\mathbf{x})]^{2}-1\}=0,
\label{eq:25}
}\end{equation}\DIFaddend 
where $c_{k}(\mathbf{x},\hat{\mathbf{p}})$ denotes the local phase speed at $\mathbf{x}$ corresponding to the $k$th sheet and with 
propagation direction $\hat{\mathbf{p}}$. At an initial point only one of these factors will vanish, say the $k$th one,
and by continuity it is the same factor that vanishes elsewhere. As a result, the eikonal equation is reduced to the simpler form
\DIFdelbegin \begin{displaymath}
\DIFdel{||\mathbf{p}(\mathbf{x})||^{2}c_{k}[\mathbf{x},\hat{\mathbf{p}}(\mathbf{x})]^{2}-1=0
}\end{displaymath}%DIFAUXCMD
\DIFdelend \DIFaddbegin \begin{equation}
\DIFadd{\|\mathbf{p}(\mathbf{x})\|^{2}c_{k}[\mathbf{x},\hat{\mathbf{p}}(\mathbf{x})]^{2}-1=0,
\label{eq:26}
}\end{equation}\DIFaddend 
and this is the equation we will aim to solve.

\subsection*{The method of characteristics}
To solve the reduced form of the eikonal equation we will proceed in a way that will likely seem rather unmotivated. We 
introduce the \textbf{ray Hamiltonian} for the $k$th sheet of the local slowness surface
\DIFdelbegin \begin{displaymath}
\DIFdel{H_{k}(\mathbf{x},\mathbf{p})=\frac{1}{2}||\mathbf{p}(\mathbf{x})||^{2}c_{k}[\mathbf{x},\hat{\mathbf{p}}(\mathbf{x})]^{2},
}\end{displaymath}%DIFAUXCMD
\DIFdelend \DIFaddbegin \begin{equation}
\DIFadd{H_{k}(\mathbf{x},\mathbf{p})=\frac{1}{2}\|\mathbf{p}\|^{2}c_{k}(\mathbf{x},\hat{\mathbf{p}})^{2},
\label{eq:27}
}\end{equation}\DIFaddend 
which is defined on a six-dimensional \textbf{phase space} \DIFdelbegin \DIFdel{comprised of }\DIFdelend \DIFaddbegin \DIFadd{comprising }\DIFaddend points $(\mathbf{x}, \mathbf{p})$, and allows us 
to write the eikonal equation\DIFdelbegin \DIFdel{eq.(26)in the alternate form
}\begin{displaymath}
\DIFdel{H_{k}[\mathbf{x},\mathbf{p}(\mathbf{x})]=\frac{1}{2}, \quad \mathbf{p}(\mathbf{x})=(\nabla T)(\mathbf{x}).
}\end{displaymath}%DIFAUXCMD
\DIFdelend \DIFaddbegin \DIFadd{, eq.~(\ref{eq:26}), in the alternative form
}\begin{equation}
\DIFadd{H_{k}[\mathbf{x},\mathbf{p}(\mathbf{x})]=\frac{1}{2}, \quad \mathbf{p}(\mathbf{x})=(\nabla T)(\mathbf{x}).
\label{eq:28}
}\end{equation}\DIFaddend 
Eq.\DIFdelbegin \DIFdel{(28}\DIFdelend \DIFaddbegin \DIFadd{~(\ref{eq:28}}\DIFaddend ) states that the desired solution, $T$, of the eikonal equation corresponds to a \textbf{three-dimensional \DIFdelbegin \DIFdel{level surfaces}\DIFdelend \DIFaddbegin \DIFadd{surface}\DIFaddend } \DIFaddbegin \DIFadd{$S=\{(\mathbf{x},\nabla T(\mathbf{x}))\}$ lying within 
 the level set }\DIFaddend $H_{k}(\mathbf{x},\mathbf{p})=\frac{1}{2}$ of the ray Hamiltonian. This 
geometric viewpoint is central to the method of characteristics. 

The \textbf{initial data} we have for the problem is that on the plane
\DIFdelbegin \begin{displaymath}
\DIFdel{p_{i}^{0}x_{i}=0, 
}\end{displaymath}%DIFAUXCMD
\DIFdelend \DIFaddbegin \begin{equation}
\DIFadd{p_{i}^{0}x_{i}=0, 
\label{eq:29}
}\end{equation}\DIFaddend 
the travel time is constant, and we are free to set this value equal to zero. We then wish to continue this solution into 
the heterogeneous region $p_{i}^{0}x_{i}\ge0$. 
Each point $\mathbf{x}^{0}$ on this initial plane corresponds to a point $(\mathbf{x}^{0},\mathbf{p}^{0})$ within the phase space, 
and at such points eq.\DIFdelbegin \DIFdel{(28}\DIFdelend \DIFaddbegin \DIFadd{~(\ref{eq:28}}\DIFaddend ) is satisfied.
In this manner, the initial data for the problem specifies a \textbf{two-dimensional \DIFdelbegin \DIFdel{level }\DIFdelend surface $S^{0}$ \DIFdelbegin \DIFdel{of the ray Hamiltonian}\DIFdelend \DIFaddbegin \DIFadd{lying within this level set}\DIFaddend }, and to 
solve the eikonal equation \DIFdelbegin \DIFdel{, }\DIFdelend we must construct
 a three-dimensional \DIFdelbegin \DIFdel{level }\DIFdelend surface $S$ \DIFdelbegin \DIFdel{containing }\DIFdelend \DIFaddbegin \DIFadd{within the level set that contains }\DIFaddend $S^{0}$.

Let $(\mathbf{x}^{0},\mathbf{p}^{0})$ be a point on the initial \DIFdelbegin \DIFdel{plane }\DIFdelend \DIFaddbegin \DIFadd{surface }\DIFaddend $S^{0}$, and consider the solution of 
\textbf{Hamilton's canonical equations} for the ray Hamiltonian $H_{k}$ as defined by
\DIFdelbegin \begin{displaymath}
\DIFdel{\frac{\ddns x_{i}}{d\sigma}=\frac{\partial H_{k}}{\partial p_{i}}, \quad \frac{\ddns p_{i}}{d\sigma}=-\frac{\partial H_{k}}{\partial x_{i}}
}\end{displaymath}%DIFAUXCMD
\DIFdelend \DIFaddbegin \begin{equation}
\DIFadd{\frac{\ddns x_{i}}{\ddns\sigma}=\frac{\partial H_{k}}{\partial p_{i}}, \quad \frac{\ddns p_{i}}{\ddns\sigma}=-\frac{\partial H_{k}}{\partial x_{i}},
\label{eq:30}
}\end{equation}\DIFaddend 
subject to the initial conditions
\DIFdelbegin \begin{displaymath}
\DIFdel{x_{i}(0)=x_{i}^{0}, \quad p_{i}(0)=p_{i}^{0}.
}\end{displaymath}%DIFAUXCMD
\DIFdelend \DIFaddbegin \begin{equation}
\DIFadd{x_{i}(0)=x_{i}^{0}, \quad p_{i}(0)=p_{i}^{0}.
\label{eq:31}
}\end{equation}\DIFaddend 
Within these equations, $\sigma$ is a \textbf{generating parameter} whose physical interpretation will become clear but 
initially is just a parameterisation for the resulting integral curves.
A key property of Hamilton's equations is that the \textbf{Hamiltonian is conserved along their solutions}. Indeed, using the chain rule we obtain
\DIFdelbegin \begin{displaymath}
\DIFdel{\frac{\ddns}{\ddns\sigma}H}[\DIFdel{\mathbf{x}(\sigma), \mathbf{p}(\sigma)}]\DIFdel{=\frac{\partial H}{\partial x_{i}}\frac{\ddns x_{i}}{\ddns\sigma}+\frac{\partial H}{\partial p_{i}}\frac{\ddns p_{i}}{d\sigma} = -\frac{\ddns p_{i}}{d\sigma}\frac{\ddns x_{i}}{\ddns\sigma}+\frac{\ddns x_{i}}{d\sigma}\frac{\ddns p_{i}}{\ddns\sigma}=0,
}\end{displaymath}%DIFAUXCMD
\DIFdelend \DIFaddbegin \begin{equation}
\DIFadd{\frac{\ddns}{\ddns\sigma}H_{k}[\mathbf{x}(\sigma), \mathbf{p}(\sigma)]=\frac{\partial H_{k}}{\partial x_{i}}\frac{\ddns x_{i}}{\ddns\sigma}+\frac{\partial H_{k}}{\partial p_{i}}\frac{\ddns p_{i}}{\ddns\sigma} = -\frac{\ddns p_{i}}{\ddns\sigma}\frac{\ddns x_{i}}{\ddns\sigma}+\frac{\ddns x_{i}}{\ddns\sigma}\frac{\ddns p_{i}}{\ddns\sigma}=0,
\label{eq:32}
}\end{equation}\DIFaddend 
where on the \DIFdelbegin \DIFdel{right hand }\DIFdelend \DIFaddbegin \DIFadd{right-hand }\DIFaddend side we have suppressed arguments for clarity. As we have $H_{k}(\mathbf{x}^{0},\mathbf{p}^{0})=\frac{1}{2}$ \DIFaddbegin \DIFadd{at }\DIFaddend each point
 on the initial surface, it follows that $H_{k}[\mathbf{x}(\sigma), \mathbf{p}(\sigma)]=\frac{1}{2}$ along the 
curve within phase space obtained by integration of eq.\DIFdelbegin \DIFdel{(30}\DIFdelend \DIFaddbegin \DIFadd{~(\ref{eq:30}}\DIFaddend ) subject to the given initial conditions. What we can then do is take in 
turn each point $(\mathbf{x}^{0},\mathbf{p}^{0})$ of the initial \DIFdelbegin \DIFdel{plane}\DIFdelend \DIFaddbegin \DIFadd{surface $S^{0}$}\DIFaddend , and integrate Hamilton's equations to produce a 
curve on which the ray Hamiltonian is equal to one half. The union of all these curves defines a smooth three-dimensional 
surface \DIFdelbegin \DIFdel{within the phase space}\DIFdelend \DIFaddbegin \DIFadd{$S$ lying within the level set $H_{k}=\frac{1}{2}$}\DIFaddend , and this gives us the desired solution of the eikonal equation in geometric form. 

In detail, this 
solution is usually defined only locally near to $S^{0}$, and we shall return to this point in a moment. Furthermore, it is not 
actually clear that such level surfaces need to be formed from the graph of $\nabla T$ for some travel time function $T$. 
Indeed, through this process, we obtain a value for $\mathbf{p}$ for each $\mathbf{x}$, and hence a vector field 
$\mathbf{p}(\mathbf{x})$. But this vector field is meant to be the gradient of a scalar function, and hence  
its first derivatives must possess the non-trivial symmetry
\DIFdelbegin \begin{displaymath}
    \DIFdel{\frac{\partial p_{i}}{\partial x_{j}}=\frac{\partial p_{j}}{\partial x_{i}}.
}\end{displaymath}%DIFAUXCMD
\DIFdelend \DIFaddbegin \begin{equation}
    \DIFadd{\frac{\partial p_{i}}{\partial x_{j}}=\frac{\partial p_{j}}{\partial x_{i}}.
\label{eq:33}
}\end{equation}\DIFaddend 
Using the \DIFdelbegin \DIFdel{method's }\DIFdelend \DIFaddbegin \DIFadd{methods }\DIFaddend of \textbf{symplectic geometry} -- a beautiful subject that is intrinsically linked to Hamilton's equations -- 
it can be shown that the necessary symmetry of $\mathbf{p}(\mathbf{x})$ does hold, and hence a valid travel time is locally 
defined. 


\begin{figure}[ht]
\centering
\DIFdelbeginFL %DIFDELCMD < \includegraphics[width=0.8\textwidth]{figures/L5F3.png}
%DIFDELCMD < %%%
\DIFdelendFL \DIFaddbeginFL \includegraphics[width=0.8\textwidth]{figures/L5F3.pdf}
\DIFaddendFL \caption{\DIFdelbeginFL \DIFdelFL{An example of ray tracing within }\DIFdelendFL \DIFaddbeginFL \DIFaddFL{Rays from an initially plane wavefront passing through }\DIFaddendFL a \DIFdelbeginFL \DIFdelFL{heterogeneous body}\DIFdelendFL \DIFaddbeginFL \DIFaddFL{smooth region of low wave speed}\DIFaddendFL , \DIFdelbeginFL \DIFdelFL{and }\DIFdelendFL \DIFaddbeginFL \DIFaddFL{obtained by
numerical integration }\DIFaddendFL of \DIFaddbeginFL \DIFaddFL{Hamilton's equations. Behind }\DIFaddendFL the \DIFdelbeginFL \DIFdelFL{formation of }\DIFdelendFL \DIFaddbeginFL \DIFaddFL{anomaly the rays cross, and their envelope forms }\DIFaddendFL a caustic\DIFdelbeginFL \DIFdelFL{due to an area of low velocity}\DIFdelendFL .}
\DIFdelbeginFL %DIFDELCMD < \label{fig:fig3}
%DIFDELCMD < %%%
\DIFdelendFL \DIFaddbeginFL \label{fig:3}
\DIFaddendFL \end{figure}

Having constructed the \DIFdelbegin \DIFdel{level }\DIFdelend surface $S$, how do we recover the travel time $T$? We first make use of \DIFdelbegin \DIFdel{eq.(7}\DIFdelend \DIFaddbegin \DIFadd{eqs.~(\ref{eq:7}}\DIFaddend ) and (\DIFdelbegin \DIFdel{30}\DIFdelend \DIFaddbegin \DIFadd{\ref{eq:30}}\DIFaddend ) to obtain
\DIFdelbegin \begin{displaymath}
\DIFdel{\frac{\ddns}{\ddns\sigma}T}[\DIFdel{\mathbf{x}(\sigma)}]\DIFdel{=\frac{\partial T}{\partial x_{i}}\frac{\ddns x_{i}}{\ddns \sigma}=p_{i}\frac{\partial H}{\partial p_{i}}
}\end{displaymath}%DIFAUXCMD
\DIFdelend \DIFaddbegin \begin{equation}
\DIFadd{\frac{\ddns}{\ddns\sigma}T}[\DIFadd{\mathbf{x}(\sigma)}]\DIFadd{=\frac{\partial T}{\partial x_{i}}\frac{\ddns x_{i}}{\ddns \sigma}=p_{i}\frac{\partial H_{k}}{\partial p_{i}}.
\label{eq:34}
}\end{equation}\DIFaddend 
To simplify the result, we note that the ray Hamiltonian $H_{k}(\mathbf{x},\mathbf{p})$ is a homogeneous function of $\mathbf{p}$ of 
degree two, meaning
\DIFdelbegin \begin{displaymath}
\DIFdel{H_{k}(\mathbf{x},\gamma \mathbf{p})=\gamma^{2}H_{k}(\mathbf{x},\mathbf{p}),
}\end{displaymath}%DIFAUXCMD
\DIFdelend \DIFaddbegin \begin{equation}
\DIFadd{H_{k}(\mathbf{x},\gamma \mathbf{p})=\gamma^{2}H_{k}(\mathbf{x},\mathbf{p}),
\label{eq:35}
}\end{equation}\DIFaddend 
for all \DIFdelbegin \DIFdel{real $\gamma$. }\DIFdelend \DIFaddbegin \DIFadd{$\gamma>0$, and so }\DIFaddend Euler's \DIFdelbegin \DIFdel{homogeneous function theorem }\DIFdelend \DIFaddbegin \DIFadd{theorem for homogeneous functions can be applied.}\DIFaddend \footnote{
    \DIFdelbegin \DIFdel{You }\DIFdelend \DIFaddbegin \DIFadd{A proof of this theorem is given in the appendix to this lecture, but you }\DIFaddend do not need to know \DIFdelbegin \DIFdel{this theorem}\DIFdelend \DIFaddbegin \DIFadd{it}\DIFaddend . In isotropic media the necessary calculation can be performed explicitly, 
    and this is the only case \DIFaddbegin \DIFadd{in which }\DIFaddend you might be asked to reproduce this derivation.} \DIFdelbegin \DIFdel{then implies }\begin{displaymath}
\DIFdel{p_{i}\frac{\partial H}{\partial p_{i}}=2H_{k},
}\end{displaymath}%DIFAUXCMD
\DIFdelend \DIFaddbegin \DIFadd{It implies that
}\begin{equation}
\DIFadd{p_{i}\frac{\partial H_{k}}{\partial p_{i}}=2H_{k},
\label{eq:36}
}\end{equation}\DIFaddend 
and from eq.\DIFdelbegin \DIFdel{(28}\DIFdelend \DIFaddbegin \DIFadd{~(\ref{eq:28}}\DIFaddend ) we conclude that
\DIFdelbegin \begin{displaymath}
\DIFdel{\frac{\ddns }{\ddns \sigma}T}[\DIFdel{\mathbf{x}(\sigma)}]\DIFdel{=1.
}\end{displaymath}%DIFAUXCMD
\DIFdelend \DIFaddbegin \begin{equation}
\DIFadd{\frac{\ddns }{\ddns \sigma}T}[\DIFadd{\mathbf{x}(\sigma)}]\DIFadd{=1.
\label{eq:37}
}\end{equation}\DIFaddend 
On the initial surface $S^{0}$ the travel time is everywhere equal to zero, and so eq.\DIFdelbegin \DIFdel{(36}\DIFdelend \DIFaddbegin \DIFadd{~(\ref{eq:37}}\DIFaddend ) has the trivial solution
\DIFdelbegin \begin{displaymath}
\DIFdel{T}[\DIFdel{\mathbf{x}(\sigma)}]\DIFdel{=\sigma
}\end{displaymath}%DIFAUXCMD
\DIFdelend \DIFaddbegin \begin{equation}
\DIFadd{T}[\DIFadd{\mathbf{x}(\sigma)}]\DIFadd{=\sigma
\label{eq:38}
}\end{equation}\DIFaddend 
along the curve starting from each initial point $(\mathbf{x}^{0},\mathbf{p}^{0})$. In this manner, we can 
identify the generating parameter with the travel time along an integral curve. 

The curves $\sigma\mapsto \mathbf{x}(\sigma)$ formed by projection of solutions of Hamilton's equations in phase space onto the
position variable are known as \textbf{rays}. In order for the travel time $T$ constructed by the above method 
to be a \textbf{single-valued function} of $\mathbf{x}$, \DIFdelbegin \DIFdel{then }\DIFdelend it is necessary that each point $\mathbf{x}$ within the heterogeneous 
region is joined to a point $\mathbf{x}^{0}$ on the initial plane by a \textbf{unique ray}. Indeed, if there were two such 
rays, then for each we could calculate a \DIFdelbegin \DIFdel{travel-time }\DIFdelend \DIFaddbegin \DIFadd{travel time }\DIFaddend through the above method, and there is no reason why these two times 
should be equal. 

\DIFdelbegin \DIFdel{Points }\DIFdelend \DIFaddbegin \DIFadd{Where the rays cross, points }\DIFaddend $\mathbf{x}$ within the heterogeneous region \DIFdelbegin \DIFdel{that can be reached by rays starting at }\DIFdelend \DIFaddbegin \DIFadd{are reached by }\DIFaddend two or more \DIFaddbegin \DIFadd{rays starting from 
different }\DIFaddend points on the initial plane\DIFdelbegin \DIFdel{are called }\DIFdelend \DIFaddbegin \DIFadd{. Such a region of }\DIFaddend \textbf{\DIFdelbegin \DIFdel{caustics}%DIFDELCMD < \MBLOCKRIGHTBRACE%%%
\DIFdel{, and at such points }\DIFdelend \DIFaddbegin \DIFadd{multipathing}} \DIFadd{is bounded by an envelope of the ray family that is known as a 
}\textbf{\DIFadd{caustic}}\DIFadd{, and on a caustic }\DIFaddend ray theory breaks down. Though we do not have 
time to look quantitatively at the \DIFdelbegin \DIFdel{the }\DIFdelend behaviour of the scalar amplitude \DIFdelbegin \DIFdel{$A^{0}(\mathbf{x})$ }\DIFdelend \DIFaddbegin \DIFadd{$A_{0}(\mathbf{x})$ }\DIFaddend introduced in eq.\DIFdelbegin \DIFdel{(24}\DIFdelend \DIFaddbegin \DIFadd{~(\ref{eq:24}}\DIFaddend ), 
it can be shown that this ray theoretical amplitude becomes \textbf{infinite \DIFdelbegin \DIFdel{at caustics}\DIFdelend \DIFaddbegin \DIFadd{on a caustic}\DIFaddend }. In reality, caustics are characterised by 
large but finite amplitudes (the name comes from the \DIFdelbegin \DIFdel{Latin for ``to burn}\DIFdelend \DIFaddbegin \DIFadd{Greek for ``burning}\DIFaddend ''), and more 
sophisticated asymptotic techniques allow them to be handled quantitatively. \DIFaddbegin \DIFadd{The formation of a caustic behind a region of low 
wave speed is shown in Fig.~\ref{fig:3}, while Fig.~\ref{fig:4} shows rays from an initially plane wavefront passing through a 
smooth random medium. In the latter case there is nothing special about the structure, and yet caustics still form once the rays 
have propagated far enough. Their occurrence is generic, and not a peculiarity of simple models.
}

\begin{figure}[ht]
\centering
\includegraphics[width=\textwidth]{figures/L5F4.pdf}
\caption{\DIFaddFL{Rays from an initially plane wavefront passing through a smooth random medium whose wave speed fluctuates by about 
five per cent (the medium is homogeneous to the left of $x\approx1$). Caustics form generically once the rays have propagated 
a few correlation lengths of the heterogeneity.}}
\label{fig:4}
\end{figure}
\DIFaddend 

Though within this lecture \DIFdelbegin \DIFdel{, }\DIFdelend we have focused on \DIFdelbegin \DIFdel{a }\DIFdelend \DIFaddbegin \DIFadd{the }\DIFaddend problem of an initially plane wave entering into a heterogeneous region, it should
 be apparent that ray theory can be applied to problems having more general initial conditions. For example, the initial data for the 
 travel time can be specified on a curve or even at a point. Furthermore, it is possible to extend the method to allow for discontinuities
in material parameters along internal boundaries so long as they are smooth relative to the wavelength of the waves considered.

\section*{What you need to know and be able to do}
\DIFdelbegin %DIFDELCMD < \begin{itemize}
%DIFDELCMD <     %%%
\DIFdelend \DIFaddbegin \begin{enumerate}
    \DIFaddend \item[(i)] The form of the ray-series expansion, and understand its underlying physical motivation.
    \item[(ii)] How to obtain the eikonal equation and transport equations by substituting the ray expansion into the equations of motion. The full derivation given in the lecture is too involved to \DIFdelbegin \DIFdel{occur within }\DIFdelend \DIFaddbegin \DIFadd{be set in }\DIFaddend an exam, but you may be asked to apply these methods within simpler situations.
    \item[(iii)] Understand how the eikonal equation can be interpreted geometrically within a six-dimensional phase space, and be able to explain how solutions can be obtained through integration of Hamilton's equations.
    \item[(iv)] Be able to solve the Hamiltonian ray equations in simple cases, and also to deduce some of their properties. Practice in this will be given in the next lecture and within the \DIFdelbegin \DIFdel{final example sheet}\DIFdelend \DIFaddbegin \DIFadd{first problem set}\DIFaddend .
    \item[(v)] What a caustic is, and why ray theory breaks down there.
\DIFdelbegin %DIFDELCMD < \end{itemize}
%DIFDELCMD < %%%
\DIFdelend \DIFaddbegin \end{enumerate}

\appendix
\section*{\DIFadd{Appendix: Euler's theorem for homogeneous functions (non-examinable)}}
\DIFadd{A function $f$ defined on $\mathbb{R}^{n}$ is said to be }\textbf{\DIFadd{homogeneous of degree $m$}} \DIFadd{if
}\begin{equation}
\DIFadd{f(\gamma\mathbf{p})=\gamma^{m}f(\mathbf{p}),
\label{eq:39}
}\end{equation}
\DIFadd{for all $\mathbf{p}$ and all $\gamma>0$. Euler's theorem states that if such a function is differentiable, then
}\begin{equation}
\DIFadd{p_{i}\frac{\partial f}{\partial p_{i}}(\mathbf{p})=mf(\mathbf{p}).
\label{eq:40}
}\end{equation}
\DIFadd{The proof is almost trivial. Differentiating both sides of eq.~(\ref{eq:39}) with respect to $\gamma$, the chain rule applied to the left-hand side gives
}\begin{equation}
\DIFadd{p_{i}\frac{\partial f}{\partial p_{i}}(\gamma\mathbf{p})=m\gamma^{m-1}f(\mathbf{p}),
\label{eq:41}
}\end{equation}
\DIFadd{and setting $\gamma=1$ yields eq.~(\ref{eq:40}). The ray Hamiltonian $H_{k}(\mathbf{x},\mathbf{p})=\frac{1}{2}\|\mathbf{p}\|^{2}c_{k}(\mathbf{x},\hat{\mathbf{p}})^{2}$ 
is homogeneous of degree two in $\mathbf{p}$ for each fixed $\mathbf{x}$, because the local phase speed $c_{k}$ depends only on the 
direction $\hat{\mathbf{p}}$ and not on the length of $\mathbf{p}$, and so eq.~(\ref{eq:36}) follows. In an isotropic medium the result can be checked directly: 
$H_{k}=\frac{1}{2}c(\mathbf{x})^{2}\|\mathbf{p}\|^{2}$, and hence $p_{i}\,\partial H_{k}/\partial p_{i}=c(\mathbf{x})^{2}\|\mathbf{p}\|^{2}=2H_{k}$.
}\DIFaddend 

\end{document}
